\documentclass[12pt,letterpaper]{article}
\usepackage{graphicx,textcomp}
\usepackage{natbib}
\usepackage{setspace}
\usepackage{fullpage}
\usepackage{color}
\usepackage[reqno]{amsmath}
\usepackage{amsthm}
\usepackage{fancyvrb}
\usepackage{amssymb,enumerate}
\usepackage[all]{xy}
\usepackage{endnotes}
\usepackage{lscape}
\newtheorem{com}{Comment}
\usepackage{float}
\usepackage{hyperref}
\newtheorem{lem} {Lemma}
\newtheorem{prop}{Proposition}
\newtheorem{thm}{Theorem}
\newtheorem{defn}{Definition}
\newtheorem{cor}{Corollary}
\newtheorem{obs}{Observation}
\usepackage[compact]{titlesec}
\usepackage{dcolumn}
\usepackage{tikz}
\usetikzlibrary{arrows}
\usepackage{multirow}
\usepackage{xcolor}
\newcolumntype{.}{D{.}{.}{-1}}
\newcolumntype{d}[1]{D{.}{.}{#1}}
\definecolor{light-gray}{gray}{0.65}
\usepackage{url}
\usepackage{listings}
\usepackage{color}

\definecolor{codegreen}{rgb}{0,0.6,0}
\definecolor{codegray}{rgb}{0.5,0.5,0.5}
\definecolor{codepurple}{rgb}{0.58,0,0.82}
\definecolor{backcolour}{rgb}{0.95,0.95,0.92}

\lstdefinestyle{mystyle}{
	backgroundcolor=\color{backcolour},   
	commentstyle=\color{codegreen},
	keywordstyle=\color{magenta},
	numberstyle=\tiny\color{codegray},
	stringstyle=\color{codepurple},
	basicstyle=\footnotesize,
	breakatwhitespace=false,         
	breaklines=true,                 
	captionpos=b,                    
	keepspaces=true,                 
	numbers=left,                    
	numbersep=5pt,                  
	showspaces=false,                
	showstringspaces=false,
	showtabs=false,                  
	tabsize=2
}
\lstset{style=mystyle}
\newcommand{\Sref}[1]{Section~\ref{#1}}
\newtheorem{hyp}{Hypothesis}

\title{Problem Set 3}
\date{Due: March 25, 2026}
\author{Applied Stats II}


\begin{document}
	\maketitle
	\section*{Instructions}
	\begin{itemize}
	\item Please show your work! You may lose points by simply writing in the answer. If the problem requires you to execute commands in \texttt{R}, please include the code you used to get your answers. Please also include the \texttt{.R} file that contains your code. If you are not sure if work needs to be shown for a particular problem, please ask.
\item Your homework should be submitted electronically on GitHub in \texttt{.pdf} form.
\item This problem set is due before 23:59 on Wednesday March 25, 2026. No late assignments will be accepted.
	\end{itemize}

	\vspace{.25cm}
\section*{Question 1}
\vspace{.25cm}
\noindent We are interested in how governments' management of public resources impacts economic prosperity. Our data come from \href{https://www.researchgate.net/profile/Adam_Przeworski/publication/240357392_Classifying_Political_Regimes/links/0deec532194849aefa000000/Classifying-Political-Regimes.pdf}{Alvarez, Cheibub, Limongi, and Przeworski (1996)} and is labelled \texttt{gdpChange.csv} on GitHub. The dataset covers 135 countries observed between 1950 or the year of independence or the first year forwhich data on economic growth are available ("entry year"), and 1990 or the last year for which data on economic growth are available ("exit year"). The unit of analysis is a particular country during a particular year, for a total $>$ 3,500 observations. 

\begin{itemize}
	\item
	Response variable: 
	\begin{itemize}
		\item \texttt{GDPWdiff}: Difference in GDP between year $t$ and $t-1$. Possible categories include: "positive", "negative", or "no change"
	\end{itemize}
	\item
	Explanatory variables: 
	\begin{itemize}
		\item
		\texttt{REG}: 1=Democracy; 0=Non-Democracy
		\item
		\texttt{OIL}: 1=if the average ratio of fuel exports to total exports in 1984-86 exceeded 50\%; 0= otherwise
	\end{itemize}
	
\end{itemize}
\newpage
\noindent Please answer the following questions:

\begin{enumerate}
	\item Construct and interpret an unordered multinomial logit with \texttt{GDPWdiff} as the output and "no change" as the reference category, including the estimated cutoff points and coefficients.
	\item Construct and interpret an ordered multinomial logit with \texttt{GDPWdiff} as the outcome variable, including the estimated cutoff points and coefficients.
	
	
\end{enumerate}

\section*{Question 2} 
\vspace{.25cm}

\noindent Consider the data set \texttt{MexicoMuniData.csv}, which includes municipal-level information from Mexico. The outcome of interest is the number of times the winning PAN presidential candidate in 2006 (\texttt{PAN.visits.06}) visited a district leading up to the 2009 federal elections, which is a count. Our main predictor of interest is whether the district was highly contested, or whether it was not (the PAN or their opponents have electoral security) in the previous federal elections during 2000 (\texttt{competitive.district}), which is binary (1=close/swing district, 0="safe seat"). We also include \texttt{marginality.06} (a measure of poverty) and \texttt{PAN.governor.06} (a dummy for whether the state has a PAN-affiliated governor) as additional control variables. 

\begin{enumerate}
	\item [(a)]
	Run a Poisson regression because the outcome is a count variable. Is there evidence that PAN presidential candidates visit swing districts more? Provide a test statistic and p-value.

	\item [(b)]
	Interpret the \texttt{marginality.06} and \texttt{PAN.governor.06} coefficients.
	
	\item [(c)]
	Provide the estimated mean number of visits from the winning PAN presidential candidate for a hypothetical district that was competitive (\texttt{competitive.district}=1), had an average poverty level (\texttt{marginality.06} = 0), and a PAN governor (\texttt{PAN.governor.06}=1).
\end{enumerate}

\section*{Question 1: Answer}

\begin{enumerate}
	\item [(1)] Unordered Multinomial Logit
\end{enumerate}

\noindent Load the necessary packages and the data

\vspace{.25cm}

\lstinputlisting[language=R, firstline=30, lastline=30]{PS3_KC.R}

\lstinputlisting[language=R, firstline=40, lastline=40]{PS3_KC.R}

\vspace{.25cm}

\noindent Values in the column GDPWdiff were initially integers. In a new column, GDPWdiff\_category, the integers were transformed into 3 factors. If the integer was less than 0, then it was transformed to be the string "negative". If the integer was greater than 0, then it was transformed to be the string "positive". Otherwise, the integer equaled 0 and was transformed to be the string "no change".

\lstinputlisting[language=R, firstline=46, lastline=50]{PS3_KC.R}

\vspace{.25cm}

\noindent I re-leveled the factors, making "no change" the baseline category. This makes for a more intuitive analysis when comparing to the "negative" and "positive" categories.
\lstinputlisting[language=R, firstline=52, lastline=52]{PS3_KC.R}

\vspace{.25cm}

\noindent I ran an unordered multinomial logit, regressing the now factored GDPWdiff\_category against the variables REG and OIL (with no interaction). 

\lstinputlisting[language=R, firstline=55, lastline=56]{PS3_KC.R}

\begin{verbatim}
	Call:
	multinom(formula = GDPWdiff_category ~ REG + OIL, data = gdp_data, 
	Hess = TRUE)
	
	Coefficients:
	         (Intercept)      REG      OIL
	negative    3.805370 1.379282 4.783968
	positive    4.533759 1.769007 4.576321
	
	Std. Errors:
	         (Intercept)       REG      OIL
	negative   0.2706832 0.7686958 6.885366
	positive   0.2692006 0.7670366 6.885097
	
	Residual Deviance: 4678.77 
	AIC: 4690.77 
\end{verbatim}

\noindent The negative category intercept, 3.805370, represents the log-odds that a country's GDP difference will be negative, compared to the baseline category, no change, when the country is not a democracy and the country's average ratio of fuel exports to total exports in 1984-1986 did not exceed 50 percent.

\vspace{.25cm}

\noindent The positive category intercept, 4.533759, represents the log-odds that a country's GDP difference will be positive, compared to the baseline category, no change, when the country is not a democracy and the country's average ratio of fuel exports to total exports in 1984-1986 did not exceed 50 percent.

\vspace{.25cm}

\noindent The negative category REG coefficient, 1.38, represents the log-odds that a country's GDP difference will be negative (compared to the baseline category of a no change) if a country is a democracy instead of a non-democracy.

\vspace{.25cm}

\noindent The positive category REG coefficient, 1.79, represents the log-odds that a country's GDP difference will be positive (compared to the baseline category of a no change) if a country is a democracy instead of a non-democracy.

\vspace{.25cm}

\noindent The negative category OIL coefficient, 4.78, represents the log-odds that a country's GDP difference will be negative (compared to the baseline category of a no change) if a country's average ratio of fuel exports to total exports in 1984-1986 exceeded 50 percent (compared to if a country's fuel export ratio did not exceed 50 percent).

\vspace{.25cm}

\noindent The positive category OIL coefficient, 4.57, represents the log-odds that a country's GDP difference will be positive (compared to the baseline category of a no change) if a country's average ratio of fuel exports to total exports in 1984-1986 exceeded 50 percent (compared to if a country's fuel export ratio did not exceed 50 percent).

\vspace{.25cm}

\noindent I exponentiated the coefficients to find the odds, which are more interpretable than the log-odds.
\lstinputlisting[language=R, firstline=58, lastline=58]{PS3_KC.R}

\begin{verbatim}
         (Intercept)      REG       OIL
negative    44.94186 3.972047 119.57794
positive    93.10789 5.865024  97.15632
\end{verbatim}

\noindent The negative category intercept, 44.94, represents the odds (relative to the baseline category of no change) that a country's GDP difference will be negative, when the country is not a democracy and the country's average ratio of fuel exports to total exports in 1984-1986 did not exceed 50 percent.

\vspace{.25cm}

\noindent The positive category intercept, 93.11, represents the odds (relative to the baseline category of no change) that a country's GDP difference will be positive, when the country is not a democracy and the country's average ratio of fuel exports to total exports in 1984-1986 did not exceed 50 percent.

\vspace{.25cm}

\noindent In a given country, there is 3.97 times higher odds (relative to the baseline category of no change) that the GDP difference will be negative when the country is a democracy.

\vspace{.25cm}

\noindent In a given country, there is 5.87 times higher odds (relative to the baseline category of no change) that the GDP difference will be positive when the country is a democracy.

\vspace{.25cm}

\noindent In a given country, there is 119.58 times higher odds (relative to the baseline category of no change) that the GDP difference will be negative when the country's average ratio of fuel exports to total exports in 1984-1986 exceeded 50 percent.

\vspace{.25cm}

\noindent In a given country, there is 97.17 times higher odds (relative to the baseline category of no change) that the GDP difference will be positive when the country's average ratio of fuel exports to total exports in 1984-1986 exceeded 50 percent.

\begin{enumerate}
	\item [(2)] Ordered Multinomial Logit
\end{enumerate}

\noindent I created a new column, GDPWdiff\_ordered, which, similarly to the process above, was a factor transformation. The only difference is that now the factors are ordered.

\lstinputlisting[language=R, firstline=64, lastline=64]{PS3_KC.R}

\noindent I ran an ordered multinomial logit, like above, regressing GDPWdiff\_category against the variables REG and OIL (with no interaction). 

\lstinputlisting[language=R, firstline=66, lastline=67]{PS3_KC.R}

\begin{verbatim}
Call:
polr(formula = GDPWdiff_ordered ~ REG + OIL, data = gdp_data, 
Hess = TRUE)

Coefficients:
     Value Std. Error t value
REG  0.3985    0.07518   5.300
OIL -0.1987    0.11572  -1.717

Intercepts:
                     Value    Std. Error t value 
negative|no change  -0.7312   0.0476   -15.3597
no change|positive  -0.7105   0.0475   -14.9554

Residual Deviance: 4687.689 
AIC: 4695.689 
\end{verbatim}

\noindent REG: Holding other variables constant, on average, if a country is a democracy, there is a 0.3985 increase in the log odds of being in a higher GDP difference category.

\vspace{.25cm}

\noindent OIL: Holding other variables constant, on average, if a country's average ratio of fuel exports to total exports in 1984-1986 exceeded 50 percent, there is a 0.1987 decrease in the log odds of being in a higher GDP difference category.

\vspace{.25cm}

\noindent Intercept interpretations:
\noindent The cutoff between the negative and no change categories is -0.7312

\noindent The cutoff between the no change and positive category is -0.7105 

\noindent The cutoffs define the boundaries between how the model predicts category assignments. However, they are on a latent scale, so they lack interpretable units. 

\lstinputlisting[language=R, firstline=69, lastline=69]{PS3_KC.R}

\begin{verbatim}
      REG       OIL 
1.4895639 0.8197813 
\end{verbatim}

\noindent REG: Holding other variables constant, on average, if a country is a democracy, the odds of being in a higher GDP difference category increase by about 49 percent.

\vspace{.25cm}

\noindent OIL: Holding other variables constant, on average, if a country's average ratio of fuel exports to total exports in 1984-1986 exceeded 50 percent, the odds of being in a higher GDP difference category decrease by about 18 percent.

\section*{Question 2: Answer}

\noindent Note: As the instructions say to run a Poisson regression, that is what I did. However, an examination of the data reveals a disproportionate amount of zeros. So, in a real-world analysis, a Zero-inflated Poisson (ZIP) may be better for this data.

\begin{enumerate}
	\item [(a)] Run a Poisson regression because the outcome is a count variable. Is there evidence that PAN presidential candidates visit swing districts more? Provide a test statistic and p-value.
\end{enumerate}

\noindent Load the data

\lstinputlisting[language=R, firstline=89, lastline=89]{PS3_KC.R}

\noindent As this is count data, I ran a Poisson regression, regressing the number of visits by the winning candidate on whether the district was competitive, the level of poverty, and whether the state has a PAN-affiliated governor. 

\lstinputlisting[language=R, firstline=92, lastline=96]{PS3_KC.R}

\begin{verbatim}
Call:
glm(formula = PAN.visits.06 ~ competitive.district + marginality.06 + 
PAN.governor.06, family = poisson(link = "log"), data = mexico_elections)

Coefficients:
                      Estimate Std. Error z value Pr(>|z|)    
(Intercept)          -3.81023    0.22209 -17.156   <2e-16 ***
competitive.district -0.08135    0.17069  -0.477   0.6336    
marginality.06       -2.08014    0.11734 -17.728   <2e-16 ***
PAN.governor.06      -0.31158    0.16673  -1.869   0.0617 .  
---
Signif. codes:  0 ‘***’ 0.001 ‘**’ 0.01 ‘*’ 0.05 ‘.’ 0.1 ‘ ’ 1

(Dispersion parameter for poisson family taken to be 1)

Null deviance: 1473.87  on 2406  degrees of freedom
Residual deviance:  991.25  on 2403  degrees of freedom
AIC: 1299.2

Number of Fisher Scoring iterations: 7
\end{verbatim}

\noindent competitive.district has a z-value of -0.477 and a p-value of 0.6336. The p-value is greater than the critical value of 0.05. Therefore, we fail to reject the null hypothesis that PAN presidential candidates visit swing and not swing districts equally.

\begin{enumerate}
	\item [(b)] Interpret the marginality.06 and PAN.governor.06 coefficients.
\end{enumerate}

\noindent For a one-unit increase in marginality.06, holding other variables constant, on average the log count of district visits decreases by approximately 2.08.

\lstinputlisting[language=R, firstline=136, lastline=136]{PS3_KC.R}

\begin{verbatim}
[1] 0.1249302
\end{verbatim}

\noindent There are about 87 percent fewer visits to districts with higher poverty, on average and holding all else constant.

\vspace{.25cm}

\noindent For a one-unit increase in PAN.governor.06, holding other variables constant, on average the log count of district visits decreases by approximately 0.31.

\lstinputlisting[language=R, firstline=138, lastline=138]{PS3_KC.R}

\begin{verbatim}
[1] 0.733447
\end{verbatim}

\noindent There are about 27 percent fewer visits to districts in states that have PAN-affiliated governors, on average and holding all else constant.

\begin{enumerate}
	\item [(c)] Provide the estimated mean number of visits from the winning PAN presidential candidate for a hypothetical district that was competitive (\texttt{competitive.district}=1), had an average poverty level (\texttt{marginality.06} = 0), and a PAN governor (\texttt{PAN.governor.06}=1).
\end{enumerate}

\lstinputlisting[language=R, firstline=97, lastline=100]{PS3_KC.R}

\begin{verbatim}
         1 
0.01494818 
\end{verbatim}

\noindent The estimated mean number of visits from the winning PAN presidential candidate for a hypothetical district that was competitive, had an average poverty level, and a PAN governor is approximately 0.015.

\end{document}
